\subsubsection{Phân tích khám phá}

\paragraph{Thống kê mô tả}

\tablename~\ref{tab:covtype-desc-part1}, \ref{tab:covtype-desc-part2},
\ref{tab:covtype-desc-part3}, và~\ref{tab:covtype-desc-part4} cho ta thấy thống
kê mô tả các biến định lượng của bộ dữ liệu Forest CoverType.

\begin{table}[htbp]
  \centering
  \caption{Thống kê mô tả các biến định lượng của bộ dữ liệu Forest CoverType
  (phần 1)}
  \label{tab:covtype-desc-part1}
  \import{\tablespath}{covtype-desc-part1}
\end{table}

\begin{table}[htbp]
  \centering
  \caption{Thống kê mô tả các biến định lượng của bộ dữ liệu Forest CoverType
  (phần 2)}
  \label{tab:covtype-desc-part2}
  \import{\tablespath}{covtype-desc-part2}
\end{table}

\begin{table}[htbp]
  \centering
  \caption{Thống kê mô tả các biến định lượng của bộ dữ liệu Forest CoverType
  (phần 3)}
  \label{tab:covtype-desc-part3}
  \import{\tablespath}{covtype-desc-part3}
\end{table}

\begin{table}[htbp]
  \centering
  \caption{Thống kê mô tả các biến định lượng của bộ dữ liệu Forest CoverType
  (phần 4)}
  \label{tab:covtype-desc-part4}
  \import{\tablespath}{covtype-desc-part4}
\end{table}

Kết quả thống kê mô tả cho các biến định lượng trong bộ dữ liệu cho thấy sự
phân bố đa dạng về thang đo và độ biến thiên giữa các đặc trưng địa hình.

Biến \textenglish{Elevation} có giá trị trung bình khoảng 2959.37 với độ lệch
chuẩn 279.98, cho thấy sự chênh lệch đáng kể về độ cao giữa các khu vực quan
sát.

Biến \textenglish{Aspect} dao động trong khoảng từ 0 đến 360, phản ánh hướng
địa hình phân bố trên toàn bộ không gian góc phương vị.

Biến \textenglish{Slope} có giá trị trung bình 14.10 và độ lệch chuẩn 7.49,
cho thấy phần lớn khu vực có độ dốc thấp đến trung bình, tuy nhiên vẫn tồn tại
các khu vực có độ dốc lớn.

Khoảng cách đến nguồn nước gần nhất
(\textenglish{Horizontal\_Distance\_To\_Hydrology}) có giá trị trung bình
269.43, trong khi khoảng cách theo phương thẳng đứng có trung bình 46.42 và
dao động từ -173 đến 601, phản ánh sự biến thiên địa hình theo độ cao tương
đối.

Khoảng cách đến đường giao thông gần nhất
(\textenglish{Horizontal\_Distance\_To\_Roadways}) có giá trị trung bình
2350.15 với độ lệch chuẩn lớn (1559.25), cho thấy mức độ phân tán không gian
đáng kể giữa các khu vực rừng và hệ thống giao thông.

Các biến \textenglish{Hillshade\_9am}, \textenglish{Hillshade\_Noon} và
\textenglish{Hillshade\_3pm} nằm trong khoảng từ 0 đến 254, thể hiện mức độ
chiếu sáng tại các thời điểm khác nhau trong ngày. Giá trị trung bình lần lượt
là 212.15, 223.32 và 142.53, cho thấy cường độ ánh sáng cao nhất vào buổi trưa.

Biến \textenglish{Horizontal\_Distance\_To\_Fire\_Points} có giá trị trung bình
1980.29 với độ lệch chuẩn 1324.20, phản ánh sự khác biệt lớn về khoảng cách
đến các điểm cháy rừng giữa các khu vực.

Tổng thể, các biến định lượng trong bộ dữ liệu thể hiện sự đa dạng về phạm vi
giá trị và mức độ biến thiên, phản ánh đặc trưng địa hình và sinh thái phức tạp
của khu vực nghiên cứu.

\paragraph{Phân phối biến mục tiêu}

\subsubsection{Phân tích phân bố biến mục tiêu}

\tablename~\ref{tab:covtype-target-count} và
\figurename~\ref{fig:covtype-target-count-plot} cho thấy được phân phối của
biến mục tiêu trong bộ dữ liệu Forest CoverType.

\begin{table}[htbp]
  \centering
  \caption{Thống kê số lượng từng lớp của biến mục tiêu trong bộ dữ liệu Forest
  CoverType}
  \label{tab:covtype-target-count}
  \import{\tablespath}{covtype-target-count}
\end{table}

\begin{figure}[htbp]
  \centering
  \includegraphics[width=0.9\textwidth]{covtype-target-distribution}
  \caption{Biểu đồ đếm tần suất biến mục tiêu trong bộ dữ liệu Forest
  CoverType}
  \label{fig:covtype-target-count-plot}
\end{figure}

Biến mục tiêu \textenglish{Cover\_Type} trong bộ dữ liệu biểu diễn loại thảm
thực vật rừng và bao gồm 7 lớp khác nhau, tương ứng với các kiểu rừng chính
trong khu vực nghiên cứu.

Kết quả thống kê cho thấy phân bố các lớp là không cân bằng. Lớp 2 chiếm tỷ lệ
lớn nhất với 48.76\%, tiếp theo là lớp 1 với 36.46\%. Hai lớp này chiếm phần
lớn tổng số quan sát, cho thấy dữ liệu bị lệch mạnh về hai loại rừng chủ đạo.

Các lớp còn lại có tỷ lệ thấp hơn đáng kể. Cụ thể, lớp 3 chiếm 6.15\%, lớp 7
chiếm 3.53\%, lớp 6 chiếm 2.99\%, lớp 5 chiếm 1.63\%, và lớp 4 chỉ chiếm
0.47\%, là lớp hiếm nhất trong bộ dữ liệu.

Sự mất cân bằng này cho thấy bài toán phân loại rừng là một bài toán đa lớp
không cân bằng, trong đó một số lớp chiếm ưu thế rõ rệt so với các lớp còn lại.
Điều này có thể ảnh hưởng đến hiệu năng của các mô hình học máy, đặc biệt là xu
hướng thiên lệch về các lớp phổ biến nếu không có biện pháp xử lý phù hợp.

\paragraph{Các biến đặc trưng định lượng}

\figurename~\ref{fig:covtype-numeric-features-group1}
và~\ref{fig:covtype-numeric-features-group2} cho ta thấy phân phối của các biến
đặc trưng định lượng của bộ dữ liệu Forest CoverType. Hình gồm hai cột, cột
trái là biểu đồ hộp của biến đặc trưng, cột phải là biểu đồ tần suất của biến
tương ứng.

\begin{figure}[htbp]
  \centering
  \includegraphics[width=0.9\textwidth]{covtype-numeric-features-group1}
  \caption{Phân phối các biến đặc trưng định lượng của bộ dữ liệu Forest
  CoverType (phần 1)}
  \label{fig:covtype-numeric-features-group1}
\end{figure}

\begin{figure}[htbp]
  \centering
  \includegraphics[width=0.9\textwidth]{covtype-numeric-features-group2}
  \caption{Phân phối các biến đặc trưng định lượng của bộ dữ liệu Forest
  CoverType (phần 2)}
  \label{fig:covtype-numeric-features-group2}
\end{figure}

Các biến \textenglish{Slope}, \textenglish{Aspect},
\textenglish{Horizontal\_Distance\_To\_Hydrology},
\textenglish{Vertical\_Distance\_To\_Hydrology},
\textenglish{Horizontal\_Distance\_To\_Roadways} và
\textenglish{Horizontal\_Distance\_To\_Fire\_Points} có phân phối lệch phải.
Điều này cho thấy phần lớn các quan sát nằm trong các khu vực có địa hình tương
đối ``thuận lợi'', tức là độ dốc thấp đến trung bình và ở khoảng cách không quá
xa các yếu tố địa lý như nguồn nước, đường giao thông và điểm cháy rừng. Ngược
lại, chỉ một số ít quan sát thuộc các khu vực xa biệt lập hoặc có điều kiện địa
hình cực đoan, tạo thành đuôi kéo dài về phía các giá trị lớn.

Các biến \textenglish{Hillshade\_9am} và \textenglish{Hillshade\_Noon} có phân
phối lệch trái, cho thấy phần lớn các khu vực có mức độ chiếu sáng cao vào buổi
sáng và giữa trưa, tương ứng với các bề mặt ít bị che khuất bởi địa hình hoặc
thảm thực vật. Chỉ một số ít khu vực có mức chiếu sáng thấp, thường tương ứng
với các vị trí bị che bóng hoặc địa hình phức tạp.

Các biến \textenglish{Elevation} và \textenglish{Hillshade\_3pm} có phân phối
tương đối cân bằng. Đối với \textenglish{Elevation}, sự lệch trái nhẹ cho thấy
khu vực nghiên cứu chủ yếu nằm ở độ cao trung bình đến cao, với một số ít vùng
có độ cao thấp hơn. Trong khi đó, \textenglish{Hillshade\_3pm} phân bố khá đối
xứng, phản ánh sự đa dạng tương đối đồng đều của điều kiện chiếu sáng vào buổi
chiều trong khu vực khảo sát.

\paragraph{Các biến đặc trưng định tính}

\figurename~\ref{fig:covtype-categorical-features-group1},
\ref{fig:covtype-categorical-features-group2},
\ref{fig:covtype-categorical-features-group3}
và~\ref{fig:covtype-categorical-features-group4} cho ta thấy phân phối của các
đặc trưng trong bộ dữ liệu Forest CoverType.

\begin{figure}[htbp]
  \centering
  \includegraphics[width=0.9\textwidth]{covtype-categorical-features-group1}
  \caption{Biểu đồ đếm tần suất các biến đặc trưng định tính của bộ dữ liệu
  Forest CoverType (phần 1)}
  \label{fig:covtype-categorical-features-group1}
\end{figure}

\begin{figure}[htbp]
  \centering
  \includegraphics[width=0.9\textwidth]{covtype-categorical-features-group2}
  \caption{Biểu đồ đếm tần suất các biến đặc trưng định tính của bộ dữ liệu
  Forest CoverType (phần 2)}
  \label{fig:covtype-categorical-features-group2}
\end{figure}

\begin{figure}[htbp]
  \centering
  \includegraphics[width=0.9\textwidth]{covtype-categorical-features-group3}
  \caption{Biểu đồ đếm tần suất các biến đặc trưng định tính của bộ dữ liệu
  Forest CoverType (phần 3)}
  \label{fig:covtype-categorical-features-group3}
\end{figure}

\begin{figure}[htbp]
  \centering
  \includegraphics[width=0.9\textwidth]{covtype-categorical-features-group4}
  \caption{Biểu đồ đếm tần suất các biến đặc trưng định tính của bộ dữ liệu
  Forest CoverType (phần 4)}
  \label{fig:covtype-categorical-features-group4}
\end{figure}

Ta thấy \textenglish{Wilderness\_Area\_0} và~\textenglish{Wilderness\_Area\_2}
chiếm đa số trong bộ dữ liệu. Điều này có thể do độ cao trung bình của hai khu
vực này thấp (được nói ở phần giới thiệu) có có diện tích phủ lớn hơn.

Các loại đất \textenglish{Soil\_Type\_28} và \textenglish{Soil\_Type\_22} có số
lượng lớn hơn hẳn các loại đất khác.

\paragraph{Phân bố từng từng đặc trưng định lượng theo từng lớp của biến mục
tiêu}

\figurename~\ref{fig:covtype-numeric-features-vs-target-group1}
và~\ref{fig:covtype-numeric-features-vs-target-group2} cho ta thấy phân bố từng
đặc trưng theo từng lớp của biến mục tiêu trong bộ dữ liệu Forest CoverType.

\begin{figure}[htbp]
  \centering
  \includegraphics[width=0.9\textwidth]{covtype-numeric-features-vs-target-group1}
  \caption{Biểu đồ hộp từng đặc trưng định lượng theo từng lớp của biến mục
  tiêu trong bộ dữ liệu Forest CoverType (phần 1)}
  \label{fig:covtype-numeric-features-vs-target-group1}
\end{figure}

\begin{figure}[htbp]
  \centering
  \includegraphics[width=0.9\textwidth]{covtype-numeric-features-vs-target-group2}
  \caption{Biểu đồ hộp từng đặc trưng định lượng theo từng lớp của biến mục
  tiêu trong bộ dữ liệu Forest CoverType (phần 2)}
  \label{fig:covtype-numeric-features-vs-target-group2}
\end{figure}

Qua 2 hình, ta thấy biến \textenglish{Elevation} có sự phân bố khác nhau tương
đối giữa các lớp. Điều này gợi ý rằng độ cao có thể là một đặc trưng quan trọng
trong việc phân biệt các loại rừng phủ.

Các đặc trưng khác thì chưa thấy sự phân biệt nào rõ ràng.
